% ============================================================
%  CHAPITRE : FONCTION EXPONENTIELLE NÉPÉRIENNE
%  Mazava Maths — Terminale S
%  Calibré sur les sujets BAC S Madagascar 2021–2023
% ============================================================

\chapter{Fonction exponentielle népérienne}

\begin{objectifs}
	\begin{itemize}
		\item Connaître la définition et les propriétés algébriques de $\e^x$.
		\item Résoudre des équations, inéquations et systèmes avec $\e^x$.
		\item Calculer des limites classiques avec $\e^x$.
		\item Étudier des fonctions construites à partir de $\e^x$.
		\item Calculer des primitives du type $u'\e^u$.
		\item Résoudre des équations différentielles du premier ordre.
	\end{itemize}
\end{objectifs}

% ----------------------------------------------------------------
\section{Définition}
% ----------------------------------------------------------------

\begin{definition}
	La fonction \textbf{exponentielle népérienne}, notée $\exp$ ou $x\mapsto\e^x$,
	est la bijection réciproque de la fonction $\ln$ sur $\R$ :
	\[
	\forall x \in \R,\; \forall y > 0 : \quad
	y = \e^x \iff x = \ln y.
	\]
	En particulier : $\e^0 = 1$, $\e^1 = \e \approx 2{,}718$ et
	$\ln(\e^x) = x$ pour tout $x\in\R$, $\e^{\ln x} = x$ pour tout $x>0$.
\end{definition}

\begin{attention}
	$\e^x > 0$ pour tout $x \in \R$ : l'exponentielle ne s'annule jamais
	et ne prend que des valeurs \textbf{strictement positives}.
\end{attention}

% ----------------------------------------------------------------
\section{Propriétés algébriques}
% ----------------------------------------------------------------

\begin{propriete}{Règles de calcul}{}
	Pour tous réels $a$ et $b$ :
	\begin{center}
		\renewcommand{\arraystretch}{2.0}
		\begin{tabular}{|c|c|}
			\hline
			\textbf{Propriété} & \textbf{Formule} \\
			\hline
			Produit & $\e^{a+b} = \e^a \cdot \e^b$ \\
			\hline
			Quotient & $\e^{a-b} = \dfrac{\e^a}{\e^b}$ \\
			\hline
			Inverse & $\e^{-a} = \dfrac{1}{\e^a}$ \\
			\hline
			Puissance & $\left(\e^a\right)^n = \e^{na}$, \quad $n\in\Z$ \\
			\hline
			Puissance réelle & $\left(\e^a\right)^\alpha = \e^{\alpha a}$, \quad $\alpha\in\R$ \\
			\hline
			Valeur en 0 & $\e^0 = 1$ \\
			\hline
		\end{tabular}
	\end{center}
\end{propriete}

\begin{attention}
	$\e^{a+b} \neq \e^a + \e^b$ \quad et \quad $\e^{a \times b} \neq \e^a \times \e^b$.
\end{attention}

\begin{exemple}
	Simplifier les expressions suivantes :
	\[
	A = \e^3 \cdot \e^{-1}, \qquad
	B = \frac{\e^{2x+1}}{\e^{x-1}}, \qquad
	C = \left(\e^{3x}\right)^2 \cdot \e^{-4x}.
	\]
	\[
	A = \e^{3-1} = \e^2.
	\]
	\[
	B = \e^{(2x+1)-(x-1)} = \e^{x+2}.
	\]
	\[
	C = \e^{6x} \cdot \e^{-4x} = \e^{2x}.
	\]
\end{exemple}

% ----------------------------------------------------------------
\section{Limites usuelles}
% ----------------------------------------------------------------

\begin{propriete}{Limites aux bornes et limites classiques}{}
	\begin{center}
		\renewcommand{\arraystretch}{2.2}
		\begin{tabular}{|c|c|}
			\hline
			\textbf{Limite} & \textbf{Résultat} \\
			\hline
			$\displaystyle\lim_{x\to-\infty}\e^x$ & $0$ \\
			\hline
			$\displaystyle\lim_{x\to+\infty}\e^x$ & $+\infty$ \\
			\hline
			$\displaystyle\lim_{x\to-\infty} x\e^x$ & $0$ \quad ($\e^x$ écrase $x$) \\
			\hline
			$\displaystyle\lim_{x\to+\infty} \dfrac{\e^x}{x^n}$ \quad ($n\geq 1$) & $+\infty$ \quad ($\e^x$ domine $x^n$) \\
			\hline
			$\displaystyle\lim_{x\to-\infty} x^n\e^x$ \quad ($n\geq 1$) & $0$ \quad ($\e^x$ écrase $x^n$) \\
			\hline
			$\displaystyle\lim_{x\to 0}\dfrac{\e^x - 1}{x}$ & $1$ \\
			\hline
		\end{tabular}
	\end{center}
\end{propriete}

\begin{astucebac}
	La limite $\displaystyle\lim_{x\to 0}\frac{\e^x-1}{x}=1$ est utilisée au BAC
	pour lever des formes indéterminées $\frac{0}{0}$
	et pour étudier la dérivabilité en $x_0=0$.
	Retenir également : $\e^x$ \textbf{l'emporte toujours} sur toute puissance de $x$.
\end{astucebac}

\begin{exemple}
	Calculer $\displaystyle L = \lim_{x\to+\infty}\frac{x^2+3}{2\e^x}$.
	
	On écrit $\dfrac{x^2+3}{2\e^x} = \dfrac{1}{2}\cdot\dfrac{x^2}{\e^x} + \dfrac{3}{2\e^x}$.
	Or $\displaystyle\lim_{x\to+\infty}\frac{x^2}{\e^x}=0$ et
	$\displaystyle\lim_{x\to+\infty}\frac{1}{\e^x}=0$, donc $L = 0$.
\end{exemple}

% ----------------------------------------------------------------
\section{Dérivée et fonctions composées}
% ----------------------------------------------------------------

\begin{propriete}{Dérivées avec $\e^x$}{}
	\begin{center}
		\renewcommand{\arraystretch}{2.2}
		\begin{tabular}{|c|c|c|}
			\hline
			\textbf{Fonction} & \textbf{Dérivée} & \textbf{Domaine} \\
			\hline
			$\e^x$ & $\e^x$ & $\R$ \\
			\hline
			$\e^{u(x)}$ & $u'(x)\cdot\e^{u(x)}$ & $\R$ \\
			\hline
			$\e^{ax+b}$ & $a\,\e^{ax+b}$ & $\R$ \\
			\hline
			$\e^{-x}$ & $-\e^{-x}$ & $\R$ \\
			\hline
		\end{tabular}
	\end{center}
\end{propriete}

\begin{methode}
	Pour dériver $\e^{u(x)}$ :
	\begin{enumerate}
		\item Identifier $u(x)$.
		\item Calculer $u'(x)$.
		\item Appliquer : $\bigl(\e^u\bigr)' = u' \cdot \e^u$.
	\end{enumerate}
\end{methode}

\begin{exemple}
	Calculer les dérivées des fonctions suivantes :
	\[
	f(x) = \e^{3x^2-1}, \qquad
	g(x) = x^2\e^{-x}, \qquad
	h(x) = \frac{\e^x - 1}{\e^x + 1}.
	\]
	\[
	f'(x) = 6x\,\e^{3x^2-1}.
	\]
	\[
	g'(x) = 2x\e^{-x} + x^2(-1)\e^{-x} = x(2-x)\e^{-x}.
	\]
	\[
	h'(x) = \frac{\e^x(\e^x+1) - \e^x(\e^x-1)}{(\e^x+1)^2}
	= \frac{2\e^x}{(\e^x+1)^2}.
	\]
\end{exemple}

% ----------------------------------------------------------------
\section{Étude complète de $x\mapsto\e^x$}
% ----------------------------------------------------------------

\begin{propriete}{Tableau de variation de $\e^x$}{}
	\begin{center}
		\renewcommand{\arraystretch}{1.6}
		\begin{tabular}{|c|ccc|}
			\hline
			$x$ & $-\infty$ & & $+\infty$ \\
			\hline
			$(\e^x)' = \e^x$ & \multicolumn{3}{c|}{$+$} \\
			\hline
			$\e^x$ & $0$ & $\nearrow$ & $+\infty$ \\
			\hline
		\end{tabular}
	\end{center}
	La fonction $\e^x$ est \textbf{strictement croissante} sur $\R$,
	positive, sans extremum.
\end{propriete}

% Figure TikZ — courbe de exp
\begin{center}
	\begin{tikzpicture}[>=Stealth, scale=1.0]
		\draw[->] (-3.2,0) -- (3.2,0) node[right,font=\small]{$x$};
		\draw[->] (0,-0.3) -- (0,5.0) node[above,font=\small]{$y$};
		% graduations x
		\foreach \x in {-3,-2,-1,1,2,3}{
			\draw(\x,2pt)--(\x,-2pt) node[below,font=\tiny]{\x};
		}
		% graduations y
		\foreach \y in {1,2,3,4}{
			\draw(2pt,\y)--(-2pt,\y) node[left,font=\tiny]{\y};
		}
		% courbe exp
		\draw[thick, domain=-3.1:1.6, samples=80]
		plot(\x,{exp(\x)});
		% tangente en x=0 (pente 1, y=x+1)
		\draw[dashed,gray,domain=-2:2]
		plot(\x,{\x+1}) node[right,font=\tiny,gray]{$y=x+1$};
		% point (0,1)
		\filldraw (0,1) circle(2pt) node[above right,font=\small]{$(0,1)$};
		% asymptote
		\draw[dashed,gray!60] (-3.1,0)--(3.1,0);
		\node[right,font=\small] at (1.5,{exp(1.5)}){$\mathcal{C}_{\exp}$};
		\node[below left,font=\small] at(0,0){$O$};
	\end{tikzpicture}
\end{center}

\begin{remarque}
	Les courbes de $\ln$ et de $\e^x$ sont \textbf{symétriques}
	par rapport à la droite $y=x$
	(ce sont des fonctions réciproques l'une de l'autre).
\end{remarque}

% ----------------------------------------------------------------
\section{Équations et inéquations avec $\e^x$}
% ----------------------------------------------------------------

\begin{methode}
	\textbf{Équations :}
	\begin{itemize}
		\item $\e^{u(x)} = a$ avec $a>0$ : $u(x) = \ln a$.
		\item $\e^{u(x)} = \e^{v(x)}$ : $u(x) = v(x)$ \quad (bijectivité).
		\item $\e^{u(x)} = 0$ : \textbf{impossible} (pas de solution).
	\end{itemize}
	\textbf{Inéquations :} $\e^x$ est croissante, donc :
	\[
	\e^{u(x)} \leq \e^{v(x)} \iff u(x) \leq v(x).
	\]
\end{methode}

\begin{exemple}
	Résoudre dans $\R$ :
	\[
	\text{a)}\quad \e^{2x-1} = \e^{x+3}
	\qquad
	\text{b)}\quad \e^{x^2} \leq \e^{2x}
	\qquad
	\text{c)}\quad \e^{2x} - 3\e^x + 2 = 0
	\]
	
	\textbf{a)} $2x-1 = x+3 \implies x = 4$. \quad $S=\{4\}$.
	
	\textbf{b)} $x^2 \leq 2x \implies x^2-2x\leq 0
	\implies x(x-2)\leq 0 \implies x\in[0\,;2]$.
	
	\textbf{c)} On pose $X = \e^x > 0$ :
	$X^2 - 3X + 2 = 0 \implies (X-1)(X-2)=0$.
	$X=1 \implies \e^x=1 \implies x=0$ ;
	$X=2 \implies \e^x=2 \implies x=\ln 2$.
	$S = \{0\,;\,\ln 2\}$.
\end{exemple}

% ----------------------------------------------------------------
\section{Primitives avec $\e^x$}
% ----------------------------------------------------------------

\begin{propriete}{Primitives faisant intervenir $\e^x$}{}
	\begin{center}
		\renewcommand{\arraystretch}{2.2}
		\begin{tabular}{|c|c|}
			\hline
			\textbf{Fonction} & \textbf{Primitive} \\
			\hline
			$\e^x$ & $\e^x + C$ \\
			\hline
			$\e^{ax+b}$ & $\dfrac{1}{a}\e^{ax+b} + C$ \\
			\hline
			$u'(x)\,\e^{u(x)}$ & $\e^{u(x)} + C$ \\
			\hline
			$\e^{-x}$ & $-\e^{-x} + C$ \\
			\hline
		\end{tabular}
	\end{center}
\end{propriete}

\begin{astucebac}
	Dès que vous voyez $u'\e^u$, la primitive est $\e^u$.
	Si le coefficient n'est pas exactement $u'$, ajustez par un facteur.
\end{astucebac}

\begin{exemple}
	Calculer :
	$\displaystyle\int_0^1 (2x+1)\e^{x^2+x}\dd x$.
	
	On pose $u = x^2+x$, $u' = 2x+1$.
	On reconnaît $u'\e^u$, donc :
	\[
	\int_0^1 (2x+1)\e^{x^2+x}\dd x
	= \Big[\e^{x^2+x}\Big]_0^1
	= \e^2 - \e^0 = \e^2 - 1.
	\]
\end{exemple}

% ----------------------------------------------------------------
\section{Fonction puissance $x\mapsto a^x$}
% ----------------------------------------------------------------

\begin{definition}
	Pour tout réel $a>0$ et tout réel $x$ :
	\[
	a^x = \e^{x\ln a}.
	\]
\end{definition}

\begin{propriete}{Dérivée de $a^x$}{}
	\[
	\left(a^x\right)' = \ln(a)\cdot a^x.
	\]
	En particulier, si $u$ est dérivable :
	$\left(a^{u(x)}\right)' = u'(x)\ln(a)\cdot a^{u(x)}$.
\end{propriete}

% ----------------------------------------------------------------
\section{Équations différentielles du premier ordre}
% ----------------------------------------------------------------

\begin{definition}
	Une \textbf{équation différentielle} est une équation dont l'inconnue
	est une \textbf{fonction} et qui fait intervenir cette fonction
	et ses dérivées.
\end{definition}

\begin{theoreme}{Solution de $y' + ay = 0$}
	L'équation différentielle $y' + ay = 0$ ($a\in\R$) a pour
	\textbf{solution générale} sur $\R$ :
	\[
	\boxed{y(x) = C\,\e^{-ax}, \quad C \in \R.}
	\]
\end{theoreme}

\begin{methode}
	Pour résoudre $y' = ky$ (ou $y' - ky = 0$) :
	\begin{enumerate}
		\item Reconnaître la forme $y' + ay = 0$ avec $a = -k$.
		\item Écrire la solution générale $y = C\e^{-ax} = C\e^{kx}$.
		\item Si une condition initiale $y(x_0)=y_0$ est donnée,
		substituer pour trouver $C$.
	\end{enumerate}
\end{methode}

\begin{exemple}
	Résoudre $y' - 2y = 0$ avec $y(0) = 3$.
	
	Solution générale : $y = C\e^{2x}$.
	Condition : $y(0) = C = 3$.
	\[
	\boxed{y(x) = 3\e^{2x}.}
	\]
\end{exemple}

% ----------------------------------------------------------------
\section{Exercices résolus}
% ----------------------------------------------------------------

\exercice{\etoilepleine\etoilevide\etoilevide}
\begin{enumerate}
	\item Simplifier : $A = \dfrac{\e^{3x}\cdot\e^{-x}}{\e^{x+1}}$.
	\item Résoudre dans $\R$ : $\e^{2x} - 5\e^x + 6 = 0$.
	\item Calculer $f'(x)$ pour $f(x) = (x^2+1)\e^{-2x}$.
\end{enumerate}

\solution

\textbf{1.}
$A = \dfrac{\e^{3x-x}}{\e^{x+1}} = \dfrac{\e^{2x}}{\e^{x+1}} = \e^{2x-(x+1)} = \e^{x-1}.$

\textbf{2.} On pose $X=\e^x>0$ : $X^2-5X+6=0 \implies (X-2)(X-3)=0$.
$X=2 \implies x=\ln 2$ ; $X=3 \implies x=\ln 3$. \quad $S=\{\ln 2\,;\,\ln 3\}$.

\textbf{3.}
$f'(x) = 2x\,\e^{-2x} + (x^2+1)(-2)\e^{-2x} = \e^{-2x}(2x - 2x^2 - 2) = -2\e^{-2x}(x^2-x+1).$

\bigskip
\exercice{\etoilepleine\etoilepleine\etoilevide}
Soit $f$ définie sur $\R$ par $f(x) = (x-1)\e^x + 2$.
\begin{enumerate}
	\item Calculer $\displaystyle\lim_{x\to-\infty}f(x)$ et
	$\displaystyle\lim_{x\to+\infty}f(x)$.
	\item Calculer $f'(x)$ et dresser le tableau de variation de $f$.
	\item Montrer que l'équation $f(x)=0$ admet une unique solution $\alpha$
	sur $\R$, et préciser un encadrement de $\alpha$.
\end{enumerate}

\solution

\textbf{1.}
$\displaystyle\lim_{x\to-\infty}(x-1)\e^x = 0$
(car $x\e^x\to 0$), donc $\displaystyle\lim_{x\to-\infty}f(x) = 0+2 = 2$.

$\displaystyle\lim_{x\to+\infty}(x-1)\e^x = +\infty$,
donc $\displaystyle\lim_{x\to+\infty}f(x) = +\infty$.

\textbf{2.}
$f'(x) = \e^x + (x-1)\e^x = x\e^x$.

$\e^x > 0$ toujours, donc $f'(x)$ a le signe de $x$ :

\begin{center}
	\renewcommand{\arraystretch}{1.6}
	\begin{tabular}{|c|ccccc|}
		\hline
		$x$ & $-\infty$ & & $0$ & & $+\infty$ \\
		\hline
		$f'(x)$ & & $-$ & $0$ & $+$ & \\
		\hline
		$f(x)$ & $2$ & $\searrow$ & $1$ & $\nearrow$ & $+\infty$ \\
		\hline
	\end{tabular}
\end{center}

\textbf{3.} $f$ admet un minimum de $1>0$ en $x=0$.
Or $\displaystyle\lim_{x\to-\infty}f(x)=2>0$ et $f$ est toujours $\geq 1$.
\textbf{Attention :} $f(x)\geq 1 > 0$ pour tout $x$, donc $f(x)=0$
n'admet \textbf{aucune solution} sur $\R$.

\begin{remarque}
	Cet exercice illustre l'importance du minimum : avant d'invoquer le TVI,
	vérifier que la fonction change de signe.
\end{remarque}

\bigskip
\exercice{\etoilepleine\etoilepleine\etoilepleine}
Soit $f$ définie sur $[0\,;+\infty[$ par :
$f(0)=-2$ et $f(x)=x-2+\e^{-\frac{1}{x}}$ pour $x>0$.

\begin{enumerate}
	\item Étudier la continuité et la dérivabilité de $f$ à droite en $x_0=0$.
	\item Calculer $f'(x)$ pour $x>0$ et dresser le tableau de variation.
	\item Montrer que $\Delta : y=x-1$ est asymptote à $\mathcal{C}$ en $+\infty$.
	Étudier la position de $\mathcal{C}$ par rapport à $\Delta$.
\end{enumerate}

\solution

\textbf{1. Continuité à droite en 0.}

Quand $x\to 0^+$ : $\dfrac{1}{x}\to+\infty$, donc $\e^{-1/x}\to 0$.
\[
\lim_{x\to 0^+}f(x) = 0-2+0 = -2 = f(0). \quad \text{Donc $f$ est continue à droite en 0.}
\]

Dérivabilité à droite :
\[
\lim_{x\to 0^+}\frac{f(x)-f(0)}{x-0}
= \lim_{x\to 0^+}\frac{x-2+\e^{-1/x}+2}{x}
= \lim_{x\to 0^+}\left(1 + \frac{\e^{-1/x}}{x}\right).
\]
Or $\dfrac{\e^{-1/x}}{x} = \dfrac{1/x}{\e^{1/x}} \to 0$ quand $x\to 0^+$
(car $\e^t/t\to+\infty$). Donc le taux est $1$ : $f$ admet une demi-tangente
de pente $1$ en $x_0=0$.

\textbf{2. Dérivée pour $x>0$.}
\[
f'(x) = 1 + \frac{1}{x^2}\e^{-1/x}.
\]
$f'(x)>0$ pour tout $x>0$ ($\e^{-1/x}>0$), donc $f$ est
\textbf{strictement croissante} sur $]0\,;+\infty[$.

\begin{center}
	\renewcommand{\arraystretch}{1.6}
	\begin{tabular}{|c|ccc|}
		\hline
		$x$ & $0$ & & $+\infty$ \\
		\hline
		$f'(x)$ & \multicolumn{3}{c|}{$+$} \\
		\hline
		$f(x)$ & $-2$ & $\nearrow$ & $+\infty$ \\
		\hline
	\end{tabular}
\end{center}

\textbf{3. Asymptote.}
\[
f(x)-(x-1) = x-2+\e^{-1/x} - x + 1 = -1+\e^{-1/x}
\xrightarrow[x\to+\infty]{} -1+0 = -1 \neq 0.
\]

Donc $\Delta: y=x-1$ n'est \textbf{pas} asymptote.
Testons $y=x-2$ :
\[
f(x)-(x-2) = \e^{-1/x} \xrightarrow[x\to+\infty]{} 0.
\]
Donc $\Delta: y=x-2$ est l'asymptote oblique de $\mathcal{C}$ en $+\infty$.

Position : $f(x)-(x-2) = \e^{-1/x} > 0$ pour $x>0$,
donc $\mathcal{C}$ est \textbf{au-dessus} de $\Delta$.

% ----------------------------------------------------------------
\section{Exercice type BAC \textnormal{(d'après sujet 2022 -- Problème 2)}}
% ----------------------------------------------------------------

\exercice{\etoilepleine\etoilepleine\etoilepleine}

Soit $f$ la fonction définie sur $\R$ par :
\[
f_n(x) = \e^{nx}\ln(1+\e^{-nx}), \quad n\in\N^*.
\]
On désigne par $(\mathcal{C}_n)$ sa courbe représentative dans un repère
orthonormé d'unité $2$\,cm.

\medskip
\textbf{Partie A : Étude de $f_n$}
\begin{enumerate}
	\item Calculer $\displaystyle\lim_{x\to-\infty}f_n(x)$ et
	$\displaystyle\lim_{x\to+\infty}f_n(x)$ (poser $X=\e^{-nx}$).
	\item On considère $g_n(x) = \ln(1+\e^{-nx}) - \dfrac{1}{1+\e^{nx}}$.
	\begin{enumerate}
		\item Vérifier que $g_n$ est strictement décroissante sur $\R$.
		\item En déduire que $g_n(x)>0$ pour tout $x\in\R$.
	\end{enumerate}
	\item Montrer que $f_n'(x) = n\e^{nx}\cdot g_n(x)$, puis dresser
	le tableau de variation de $f_n$.
	\item Construire $(\mathcal{C}_2)$.
\end{enumerate}

\textbf{Partie B : Équation différentielle}

On considère l'équation différentielle $(E)$ :
\[
y' - 2y = \frac{-2}{1+\e^{-2x}}.
\]
\begin{enumerate}
	\item Vérifier que $f_2(x)=\e^{2x}\ln(1+\e^{-2x})$ est solution de $(E)$.
	\item Résoudre dans $\R$ l'équation homogène $(E')$ : $y'-2y=0$.
	\item Montrer qu'une fonction $\varphi$ est solution de $(E)$
	si et seulement si $\varphi - f_2$ est solution de $(E')$.
	\item En déduire la solution générale de $(E)$.
\end{enumerate}

\solution

\textbf{Partie A.}

\textbf{1.} Posons $X=\e^{-nx}$.
Quand $x\to-\infty$ : $X\to+\infty$ et $\e^{nx}=1/X\to 0$,
donc $f_n(x) = \dfrac{\ln(1+X)}{X}\cdot X\cdot\e^{nx}\to 0\cdot +\infty$...
Plus précisément : $f_n(x)=\dfrac{\ln(1+X)}{X}$ avec $X\to+\infty$,
donc $f_n\to 0$ (ln négligeable devant $X$). Ainsi $\displaystyle\lim_{x\to-\infty}f_n(x)=0$.

Quand $x\to+\infty$ : $X=\e^{-nx}\to 0^+$, $\ln(1+X)\sim X$,
donc $f_n(x)\sim \e^{nx}\cdot\e^{-nx}=1$... Plus précisément :
$\ln(1+X)\to\ln 1=0$ et $\e^{nx}\to+\infty$, donc forme $\infty\cdot 0$.
On écrit $f_n(x)=\dfrac{\ln(1+\e^{-nx})}{\e^{-nx}}$; quand $x\to+\infty$
et $X\to 0$ : $\dfrac{\ln(1+X)}{X}\to 1$, donc $f_n(x)\to 1$.
Ainsi $\displaystyle\lim_{x\to+\infty}f_n(x)=1$ (en utilisant $\lim_{X\to 0}\frac{\ln(1+X)}{X}=1$).

\textbf{2a.}
$g_n'(x) = \dfrac{-n\e^{-nx}}{1+\e^{-nx}} + \dfrac{n\e^{nx}}{(1+\e^{nx})^2}$.

On peut vérifier (après calcul) que $g_n'(x)<0$ pour tout $x$,
donc $g_n$ est strictement décroissante.

\textbf{2b.} $g_n(0)=\ln 2 - \dfrac{1}{2} > 0$ (car $\ln 2\approx 0{,}693>0{,}5$).
Comme $g_n$ est décroissante et $\displaystyle\lim_{x\to+\infty}g_n(x)=0^+$,
on a $g_n(x)>0$ pour tout $x\in\R$.

\textbf{3.} $f_n'(x)=n\e^{nx}g_n(x)>0$ pour tout $x$ (car $\e^{nx}>0$ et $g_n>0$).
Donc $f_n$ est \textbf{strictement croissante} sur $\R$.

\begin{center}
	\renewcommand{\arraystretch}{1.6}
	\begin{tabular}{|c|ccc|}
		\hline
		$x$ & $-\infty$ & & $+\infty$ \\
		\hline
		$f_n'$ & \multicolumn{3}{c|}{$+$} \\
		\hline
		$f_n$ & $0$ & $\nearrow$ & $1$ \\
		\hline
	\end{tabular}
\end{center}

\textbf{Partie B.}

\textbf{1. Vérification.} $f_2(x)=\e^{2x}\ln(1+\e^{-2x})$.
\[
f_2'(x) = 2\e^{2x}\ln(1+\e^{-2x}) + \e^{2x}\cdot\frac{-2\e^{-2x}}{1+\e^{-2x}}
= 2f_2(x) - \frac{2}{1+\e^{2x}}
= 2f_2(x) + \frac{-2}{1+\e^{2x}}.
\]
Or $\dfrac{-2}{1+\e^{2x}} = \dfrac{-2}{1+\e^{2x}}$
et $(E)$ s'écrit $y'-2y=\dfrac{-2}{1+\e^{2x}}$ (en multipliant numérateur et dénominateur de $\dfrac{-2}{1+\e^{-2x}}$ par $\e^{2x}$).
Donc $f_2'-2f_2 = \dfrac{-2}{1+\e^{2x}}$ : $f_2$ est bien solution de $(E)$. \checkmark

\textbf{2.} $(E')$ : $y'-2y=0$ : solution générale $y=C\e^{2x}$, $C\in\R$.

\textbf{3.} Soit $\varphi$ solution de $(E)$ et $w=\varphi-f_2$. Alors :
$w'-2w = (\varphi'-2\varphi)-(f_2'-2f_2) = \dfrac{-2}{1+\e^{2x}}-\dfrac{-2}{1+\e^{2x}} = 0$.
Donc $w$ est solution de $(E')$. La réciproque est immédiate.

\textbf{4.} $\varphi - f_2 = C\e^{2x}$, donc la solution générale de $(E)$ est :
\[
\boxed{\varphi(x) = C\e^{2x} + \e^{2x}\ln(1+\e^{-2x}), \quad C\in\R.}
\]